

\begin{definition}[FHE]
A FHE (FHE) scheme is a tuple (KeyGen, Encrypt, Eval, Decrypt) of PPT algorithms:
\begin{itemize}
    \item $KeyGen(1^\lambda,f) \rightarrow (pk,sk, ek)$
    \item $Encrypt_{\text{key}}(x, \tau_x) \rightarrow (c_x)$ where key is either $pk$ or $sk$
    \item $Eval_{pk}(c_x) \rightarrow (c_y)$ where $y = f(x)$
    \item $Decrypt_{key}(c_y) \rightarrow y$
    \item $Decode(y) \rightarrow \tau_y$
\end{itemize}
\end{definition}

\begin{definition}[Dataseal Verification]
A dataseal verifiable FHE (vFHE) scheme is a tuple (KeyGen, Encode Decode, Verify) of PPT algorithms:
\begin{itemize}
    \item $KeyGen(1^\lambda,f) \rightarrow (pk,sk, ek)$
    \item $Encode_{\text{key}}(x) \rightarrow (\tau_x)$ where the key is a secret encoding key (ek)
    \item $Decode(y) \rightarrow \tau_y$
    \item $Verify_{ek}(y, \tau_y) \rightarrow b$, the client accepts if $b = 1$ and rejects otherwise.
\end{itemize}
\end{definition}

\begin{definition}[Correctness]
The scheme is correct if honest computation yields correct decrypted results.
\begin{equation}
    \Pr\left[
    Dec_{sk}(c_y) = f(x)
    \;\middle|\;
    \begin{aligned}
        &(pk, sk) \leftarrow KGen(1^\lambda,f), \\
        &(c_x, \tau_x) \leftarrow Enc_{\text{key}}(x), \\
        &(c_y, \tau_y) \leftarrow Eval_{\text{key}}(c_x)
    \end{aligned}
    \right] = 1
\end{equation}
\end{definition}

\begin{definition}[Completeness]
The scheme is complete if the verify computation always accept an honest computation result.
\begin{equation}
    \Pr\left[
    \text{Verify}_{sk}(c_y, \tau_x, \tau_y) = 1
    \;\middle|\;
    \begin{aligned}
        &(pk, sk) \leftarrow \text{KGen}(1^\lambda,f), \\
        &(c_x, \tau_x) \leftarrow \text{Enc}_{\text{key}}(x), \\
        &(c_y, \tau_y) \leftarrow \text{Eval}_{\text{key}}(c_x)
    \end{aligned}
    \right] = 1
\end{equation}
\end{definition}

\begin{definition} [Soundness]
The scheme is sound if adversary cannot make the Verifier accept an incorrect or altered computation.
\begin{equation}
\Pr\left[
\begin{array}{c}
    \text{Verify}_{sk}(c_y, \tau_x, \tau_y) = 1 \\
    \land \\
    \text{Dec}_{sk}(c_y) \neq f(x)
\end{array}
\;\middle|\;
\begin{aligned}
    &(pk, sk) \leftarrow \text{KGen}(1^\lambda,f), \\
    &x \leftarrow A^{\text{OEnc},\text{ODec}}(pk), \\
    &(c_x, \tau_x) \leftarrow \text{Enc}_{\text{key}}(x), \\
    &(c_y, \tau_y) \leftarrow A^{\text{OEnc},\text{ODec}}(c_x, \tau_x)
\end{aligned}
\right] 
\end{equation}
\end{definition}

\begin{definition}[Security]
Formally, a scheme is secure if for any Probabilistic Polynomial Time (PPT) adversary $\mathcal{A}$ and any function $f$, the advantage $\text{Adv}^{\text{IND-CCA1}}_{\mathcal{A}}(\lambda) = 2\Pr[b = \hat{b}] - 1$ of the attacker in the following game is negligible in the security parameter $\lambda$:
\begin{center}
\textbf{IND-CCA1 for vFHE}
\end{center}
\begin{enumerate}
    \item $(pk, sk) \leftarrow \text{KGen}(1^\lambda, f)$
    \item $(m_0, m_1) \leftarrow \mathcal A^{{O}_{\text{Enc},\text{Dec}}}(1^\lambda, f, pk)$
    \item $(c^*, \tau^*) \leftarrow \text{Enc}_{pk}(m_b)$
    \item $\hat{b} \leftarrow \mathcal A^{{O}_{\text{Enc}}}(c^*)$
\end{enumerate}
\end{definition}